% 第2.6节最后五题，原书印刷页86--87，PDF页95--96。
\thermophaseproblem{215}{86}
把一块质量 $m=90\,\mathrm g$、
初温 $t_1=90^\circ\mathrm C$ 的铜块，
放进含有 $50\,\mathrm g$、
温度 $-3^\circ\mathrm C$ 冰的量热器中。
求达到热平衡时\emph{铜块自身}的熵变。
\exsolution
铜块从 $90^\circ\mathrm C$ 冷却释放热量。
其最多冷却至 $0^\circ\mathrm C$ 可放出
$m c_{\rm Cu}(90\,\mathrm K)
\simeq(90)(0.39)(90)\,\mathrm J
\simeq3.2\,\mathrm{kJ}$。
这些热量先把冰从 $-3^\circ\mathrm C$
加热到 $0^\circ\mathrm C$，
耗热约
$(50)(2.1)(3)\simeq315\,\mathrm J$；
余热只能熔化约
$(3200-315)/(334\,\mathrm{J/g})
\simeq8.6\,\mathrm g$ 冰，
少于原有50g，故末态仍有冰，
平衡温度恰为 $T_f=273.15\,\mathrm K$。
铜块的熵差
\[
\boxed{\Delta S_{\rm Cu}=
m c_{\rm Cu}\ln\frac{T_f}{T_1}}.
\]
取 $T_1=363.15\,\mathrm K$：
\[
\boxed{\Delta S_{\rm Cu}\simeq
(90)(0.39)\ln\frac{273.15}{363.15}
\simeq-10\,\mathrm{J/K}}.
\]
\excheck 本题只问铜块熵变，可为负；
铜块放热与冰吸热后的\emph{整个隔热体系}
熵增仍应非负。

\thermophaseproblem{216}{87}
初温 $0^\circ\mathrm C$ 的冰质量
$m_1=100\,\mathrm g$，
投入含质量
$m_2=100\,\mathrm g$ 水的绝热量热器。
水的初温分别为
(a) $t_2=60^\circ\mathrm C$；
(b) $t_2=94^\circ\mathrm C$。
忽略量热器热容，求平衡时系统总熵增。
\exsolution
记冰融点 $T_m=273.15\,\mathrm K$，
融化潜热 $q_m$，水比热 $c_w$。
(a) 热水冷至 $T_m$ 释放
$Q_h=m_2c_w(T_2-T_m)$，
若
$Q_h<m_1q_m$，
则只融化一部分冰，
末温是 $T_m$，
融冰量 $m_{\rm melt}=Q_h/q_m$。
熵增
\[
\boxed{\Delta S_a=
m_2c_w\ln\frac{T_m}{T_2}
+\frac{m_2c_w(T_2-T_m)}{T_m}}.
\]
用常用物性数值算得
$\Delta S_a\simeq9\,\mathrm{J/K}$，
与原书约 $9.2\,\mathrm{J/K}$ 相符。
(b) 若
$m_2c_w(T_2-T_m)>m_1q_m$，
全部冰先融化，剩余热量使混合水升温。
能量守恒：
\[
m_2c_w(T_2-T_f)=m_1q_m+
m_1c_w(T_f-T_m).
\]
所以
\[
\boxed{T_f=\frac{m_1T_m+m_2T_2-m_1q_m/c_w}
{m_1+m_2}}.
\]
总熵增
\[
\boxed{\Delta S_b=
\frac{m_1q_m}{T_m}
+m_1c_w\ln\frac{T_f}{T_m}
+m_2c_w\ln\frac{T_f}{T_2}}.
\]
用 $c_w\simeq4.18\,\mathrm{J/(g\,K)}$、
$q_m\simeq334\,\mathrm{J/g}$，
得到 $T_f\simeq280\,\mathrm K$，
$\Delta S_b\simeq20\,\mathrm{J/K}$。
\excheck 原书简答给约 $18\,\mathrm{J/K}$，
取不同的冰潜热和比热可影响取整；
本稿保留精确方程和所采用的物性常数，
避免在数值上强行匹配。

\thermophaseproblem{217}{87}
质量 $m=5.0\,\mathrm g$ 的熔融铅，
初温正好是铅熔点
$t_2=327^\circ\mathrm C$。
将其倒入装有大量
$0^\circ\mathrm C$ 冰的量热器中。
铅凝固潜热为
$q_l=22.5\,\mathrm{J/g}$，
铅比热
$c_l=0.125\,\mathrm{J/(g\,K)}$。
求达到热平衡时铅—冰系统总熵增。
\exsolution
冰量很大，平衡温度维持
$T_1=273.15\,\mathrm K$。
铅在熔点
$T_2\simeq600.15\,\mathrm K$ 凝固，
随后固态铅冷却至 $T_1$。
释放热量为
\[
Q=m q_l+mc_l(T_2-T_1),
\]
这些热量使一部分冰在 $T_1$ 融化，
冰的熵增是 $Q/T_1$。
铅的熵变含两段：
\[
\Delta S_{\rm Pb}
=-\frac{mq_l}{T_2}
+mc_l\ln\frac{T_1}{T_2}.
\]
故系统总熵增
\[
\boxed{\Delta S=
mq_l\left(\frac1{T_1}-\frac1{T_2}\right)
+mc_l\left(\frac{T_2}{T_1}-1
-\ln\frac{T_2}{T_1}\right)}.
\]
代入给定参数：
\[
\boxed{\Delta S\simeq0.48\,\mathrm{J/K}>0}.
\]
\excheck 使用的是铅的凝固潜热，
与冰的熔化潜热\emph{不同}；
大冰库只起温度 $T_1$ 的热库作用。

\thermophaseproblem{218}{87}
活塞下有水的饱和蒸汽，
通过压缩或膨胀使蒸汽始终保持饱和、
恰好处在凝结临界状态。
假设蒸汽为理想气体且液相比体积可忽略。
求沿此过程水蒸汽的摩尔热容 $C(T)$，
并计算 $t=100^\circ\mathrm C$ 的数值。
\exsolution
对一摩尔水蒸汽，
状态方程
$p_{\rm sat}V=RT$。
沿饱和曲线微分：
\[
p_{\rm sat}\frac{\mathrm dV}{\mathrm dT}
=R-\frac{RT}{p_{\rm sat}}
\frac{\mathrm dp_{\rm sat}}{\mathrm dT}.
\]
Clausius--Clapeyron 近似：
\[
\frac1{p_{\rm sat}}
\frac{\mathrm dp_{\rm sat}}{\mathrm dT}
\simeq\frac{q_vM}{RT^2}.
\]
代入：
\[
p_{\rm sat}\frac{\mathrm dV}{\mathrm dT}
=R-\frac{q_vM}{T}.
\]
由热力学第一定律
$C=C_V+p_{\rm sat}\mathrm dV/\mathrm dT$，
\[
\boxed{C=C_p-\frac{q_vM}{T}}.
\]
取 $T=373.15\,\mathrm K$、
$q_v\simeq2.26\,\mathrm{MJ/kg}$、
$M=0.018\,\mathrm{kg/mol}$、
水蒸汽 $C_p\simeq35\,\mathrm{J/(mol\,K)}$，
\[
\boxed{C\simeq-74\,\mathrm{J/(mol\,K)}}.
\]
\excheck 负值是\emph{饱和蒸汽沿特定相变边界}
的过程热容，不表示气体稳定态的 $C_p$ 为负。

\thermophaseproblem{219}{87}
一摩尔水起初在温度 $T_1$ 下与极微量
饱和蒸汽平衡，随后全部转变为温度
$T_2$ 下的饱和水蒸气。
忽略水的液相比体积，把蒸汽当作理想气体。
求系统熵增。
\exsolution
熵是状态函数。
构造一条假想可逆路径：
先把一摩尔液态水由 $T_1$
加热至 $T_2$，然后在 $T_2$ 下汽化。
若液态水的定压摩尔热容
$C_{p,\ell}$ 在温区内近似为常数，
升温熵增为
\[
\Delta S_{\rm warm}
=C_{p,\ell}\ln\frac{T_2}{T_1}.
\]
在 $T_2$ 时汽化，
一摩尔汽化吸热为 $Mq_v$，
对应熵增
$\Delta S_{\rm vap}=Mq_v/T_2$。
故
\[
\boxed{\Delta S=
C_{p,\ell}\ln\frac{T_2}{T_1}
+\frac{Mq_v}{T_2}}.
\]
\excheck 原书简答写 $C_p$，
本稿明确它在所选路径中是
\emph{液态水的定压摩尔热容}，
而不是第218题的水蒸汽摩尔热容；
如热容随温度变化，则应用
$\int_{T_1}^{T_2}C_{p,\ell}(T)\mathrm dT/T$。
